\subsection{A Benefit-Cost Assessment of Failed Renewals} \label{sec:mvpf}

Following \cite{hendren2020unified} and \cite{finkelstein2020welfare}, we conduct a benefit-cost analysis to evaluate the overall impact of the failed renewals dedicated to road maintenance. A failed renewal delivers an immediate and visible benefit to local households: the expiring levy returns about \$76 per year to the typical homeowner, a stream worth roughly \$1,520 in present value. Set against this benefit is a cost that arrives only years later. As the roads deteriorate, the median home loses about \$15,350 in value over the decade following the vote, as shown in Table \ref{tab:median_sale_amount}.

We summarize this imbalance with the Marginal Value of Public Funds (MVPF), a benefit-cost measure that policy analysts increasingly use to express the net return to society per dollar of public funds \citep{hendren2020unified}. We define MVPF as the ratio of household's willingness to pay (WTP) for the tax cut to the government's net cost of providing the tax cut. WTP represents the value that households place on the tax cut, equivalent to the net present value of the tax savings and the discounted loss in house values. Net cost to the government includes both the direct long-run tax revenue loss and the fiscal externality\footnote{Appendix Section \ref{sec:appxmvpf} provides details on the MVPF calculation.}. Using this framework, we compute an MVPF of -1.4, indicating that for every dollar of revenue loss to the government, the households incur a net loss of \$1.4, after accounting for both direct costs and fiscal externalities. This negative MVPF suggests that the road maintenance tax cut leads to a net loss for households\footnote{To put this in perspective, although \cite{hendren2020unified} do not report a sub-category for public infrastructure and focus on government spending programs instead of tax cuts, their MVPF estimates for policies targeting adults range from 0.5 to 2.}.

Why would voters choose a policy that leaves them worse off? The two sides of the ledger differ sharply in visibility and timing. The tax savings are immediate and salient, while road depreciation is gradual and does not register in home values until about the fourth year after the vote.
We interpret this as the public infrastructure analog to the fiscal illusion that \citet{ross2013fiscal} document for property reassessment.
In Section \ref{sec:mech}, we discuss some policy levers that could potentially address this fiscal illusion.

